
% ════════════════════════════════════════
% 附录（11.附录.tex）
% ════════════════════════════════════════
\newpage
\makeatletter
\let\@oldaddcontentsline\addcontentsline
\renewcommand{\addcontentsline}[3]{}

\section*{附录 A\quad 程序与结果文件说明}
\let\addcontentsline\@oldaddcontentsline
\makeatother

\renewcommand{\arraystretch}{1.2}
\begin{longtable}{>{\centering\arraybackslash}p{0.34\textwidth}>{\raggedright\arraybackslash}p{0.60\textwidth}}
  \caption{程序与结果文件说明表}
  \label{tab:appendix} \\
  \toprule
  文件 & 说明 \\
  \midrule
  \endfirsthead
  \caption{程序与结果文件说明表（续）} \\
  \toprule
  文件 & 说明 \\
  \midrule
  \endhead
  \bottomrule
  \endfoot
  支撑材料/code/\_models.py & 数据读取、小时预报到 10 分钟预报展开、线性规划建模公共模块 \\
  支撑材料/code/\_common.py & 绘图与结果导出公共模块 \\
  支撑材料/code/问题一.py & 问题一：单日计划购电线性规划求解程序 \\
  支撑材料/code/问题二.py & 问题二：全年逐日“预测—随机规划—结算”三段式求解程序 \\
  支撑材料/code/问题三.py & 问题三：滚动预报—偏差调整两阶段求解程序 \\
  支撑材料/code/问题四.py & 问题四：波动电价情形下的重算程序 \\
  支撑材料/预处理数据.csv & 由附件 1$\sim$附件 4 清洗得到的 10 分钟级建模数据 \\
  支撑材料/code/稳健性检验.py & 第六章稳健性检验程序（数据摄动检验、线性规划与动态规划算法一致性检验） \\
  支撑材料/pic/ & 各问题结果图，共 29 幅（问题一 6 幅、问题二 9 幅、问题三 7 幅、问题四 6 幅、模型检验 1 幅） \\
  支撑材料/pic/图1\_问题分析总体流程图.png & 问题分析总体流程图（非程序生成，手工绘制） \\
  支撑材料/result/result1.xlsx & 问题一计划购电策略结果文件 \\
  支撑材料/result/result2.xlsx & 问题二全年购电策略结果文件 \\
  支撑材料/result/result3.xlsx & 问题三含紧急购电的结果文件 \\
  支撑材料/result/result4-2.xlsx & 问题四波动电价下问题二结果文件 \\
  支撑材料/result/result4-3.xlsx & 问题四波动电价下问题三结果文件 \\
\end{longtable}

\newpage
\makeatletter
\let\@oldaddcontentsline\addcontentsline
\renewcommand{\addcontentsline}[3]{}
\section*{附录 B\quad 程序源代码}
\let\addcontentsline\@oldaddcontentsline
\makeatother

\lstset{basicstyle=\ttfamily\footnotesize,frame=single,breaklines=true,numbers=left,numberstyle=\tiny\color{gray}}

\subsection*{B.1\quad 公共模块 \texttt{\_models.py}（数据读取与优化建模）}
\lstinputlisting[language=Python]{支撑材料/code/_models.py}

\subsection*{B.2\quad 公共模块 \texttt{\_common.py}（绘图与结果导出）}
\lstinputlisting[language=Python]{支撑材料/code/_common.py}

\subsection*{B.3\quad 问题一求解程序 \texttt{问题一.py}}
\lstinputlisting[language=Python]{支撑材料/code/问题一.py}

\subsection*{B.4\quad 问题二求解程序 \texttt{问题二.py}}
\lstinputlisting[language=Python]{支撑材料/code/问题二.py}

\subsection*{B.5\quad 问题三求解程序 \texttt{问题三.py}}
\lstinputlisting[language=Python]{支撑材料/code/问题三.py}

\subsection*{B.6\quad 问题四求解程序 \texttt{问题四.py}}
\lstinputlisting[language=Python]{支撑材料/code/问题四.py}

\subsection*{B.7\quad 稳健性检验程序 \texttt{稳健性检验.py}}
\lstinputlisting[language=Python]{支撑材料/code/稳健性检验.py}
